\subsection{RQ3: Active Acquisition and Score Design}

\paragraph{Matched hidden-gold acquisition.}
We replay the proposed workflow on the first 200 Computer Science questions
using 50 selection seeds. The pool contains 142 \bc, 19 \mo, 8 \no, and
31 \bw examples, but these buckets and the benchmark answers are hidden until
selection finishes. Both disagreement-guided and random acquisition first
screen all 200 candidates with Nano and Mini and then request exactly 100 gold
labels. The disagreement policy ranks $\mathcal D$ before $\mathcal A$ and
randomizes within each stratum; the baseline samples uniformly from the full
pool. Thus both methods make 200 Nano calls, 200 Mini calls, and 100 human-label
requests. They differ only in which candidates receive gold.

Table~\ref{tab:agreement-acquisition} shows that the observable split changes
the labeled training mix. It recovers all 19
\mo and all 8 \no cases in every replay, while random annotation covers
53.1\% and 54.8\%, respectively. Both policies save 50\% of human labels
relative to labeling all 200 screened candidates. Neither saves answer-model
screening cost, which is deliberately held fixed.

\begin{table}[t]
  \centering
  \small
  \caption{Run-both-then-label acquisition on 200 Computer Science candidates.
  Each method obtains 100 gold labels after the same 200 paired screens; values
  average 50 selection seeds.}
  \label{tab:agreement-acquisition}
  \begin{tabular}{lrrrrr}
    \toprule
    Selection & \bc & \mo & \no & \bw & \mo recall\\
    \midrule
    Random & 70.92 & 10.08 & 4.38 & 14.62 & 53.1\%\\
    Disagreement-guided & 55.14 & 19.00 & 8.00 & 17.86 & 100.0\%\\
    \bottomrule
  \end{tabular}
\end{table}
\paragraph{Active acquisition coverage.}
At a fixed 100-label budget over five leakage-safe SemBench resamples, the
observable policy changes which outcomes become labeled before any router is
trained. Table~\ref{tab:active-acquisition} shows that it more than triples
Mini-only coverage and increases Nano-only coverage fivefold relative to random
acquisition. Selection first separates model agreements from disagreements and uses only
observable prediction-side signals for tie-breaking, never the hidden outcome
bucket. This establishes label-efficiency
for rare-outcome discovery; the current downstream direct router still
collapses toward all-Nano, so a matched inverse-cost WCE rerun remains
necessary.

\begin{table}[t]
  \centering
  \small
  \caption{Mean outcome composition under a 100-label acquisition budget
  across five leakage-safe SemBench resamples.}
  \label{tab:active-acquisition}
  \begin{tabular}{lrrrr}
    \toprule
    Acquisition & \bc & \mo & \no & \bw\\
    \midrule
    Random & 88.4 & 6.2 & 0.8 & 4.6\\
    Disagreement-guided & 67.2 & 21.0 & 4.0 & 7.8\\
    \bottomrule
  \end{tabular}
\end{table}

To isolate the downstream value of critical-case coverage, we also compare
random historical labels with an oracle outcome-controlled mixture. The latter
uses revealed buckets and is a diagnostic training composition, not an
implementable acquisition rule. With 100 random labels, the plain eight-token
router reaches 91.06\%
accuracy and 53.33\% critical recall. The 50/46/3/1 enriched mixture reaches
93.50\% and 93.33\% critical recall. The residual architecture narrows but does
not remove the gap: random training reaches 93.09\% accuracy and 86.67\%
critical recall, while enrichment reaches 93.50\% and 93.33\%. Enrichment uses
more Mini calls because the optimization finally sees enough recoverable
failures to recognize them.

Within enriched data, the \bw count matters. With a four-token prompt, the
50/46/3/1 mixture yields 93.09\% accuracy, 27.24\% saving, and 86.67\% critical
recall. Holding 46 \mo and 3 \no examples fixed but increasing \bw from 1 to 5
drops accuracy to 91.46\% and critical recall to 46.67\%. \bw is a noisy
conservative target: neither answer model is correct, so text features that
predict it need not align with recoverable difficulty.

\begin{table}[t]
  \centering
  \small
  \caption{Natural versus outcome-enriched 100-label training on the historical
  Movie test.}
  \label{tab:sampling}
  \begin{tabular}{lrrr}
    \toprule
    Router / sampling & Accuracy & Saving & Critical recall\\
    \midrule
    Plain / random & 91.06\% & 42.68\% & 53.33\%\\
    Plain / enriched & 93.50\% & 15.85\% & 93.33\%\\
    Residual / random & 93.09\% & 23.58\% & 86.67\%\\
    Residual / enriched & 93.50\% & 17.07\% & 93.33\%\\
    \bottomrule
  \end{tabular}
\end{table}

\subsection{RQ4: Prompt Capacity and Label Budget}

We cross three label budgets with four prompt lengths while fixing all other
settings (Table~\ref{tab:size-grid}). The best prompt length changes with
sample size: eight tokens lead at 50 and 100 labels, whereas sixteen lead at
200. More capacity is not monotonically better. At 100 labels, moving from 8 to
16 tokens drops accuracy from 93.50\% to 91.46\%, raises apparent saving to
51.22\%, and reduces critical recall from 93.33\% to 60.00\%. This is not a
benign cost trade-off; it is under-escalation.

Longer prompts also change the number of update opportunities per parameter.
We therefore match approximately 200 optimizer steps for the 16-token, 50- and
100-label settings. The 50-label result is unchanged at 92.28\% accuracy. The
100-label result improves from 91.46\% to 92.28\% and critical recall from
60.00\% to 73.33\%, but remains below the eight-token configuration. Update
count explains part, but not all, of the interaction between prompt capacity and
sample size.

\begin{table}[t]
  \centering
  \small
  \caption{Train size $\times$ soft-prompt length on historical Movie data.
  Each cell is mean answer accuracy / Mini saving (\%).}
  \label{tab:size-grid}
  \begin{tabular}{lrrrr}
    \toprule
    Labels & 2 tok & 4 tok & 8 tok & 16 tok\\
    \midrule
    50  & 92.28/50.41 & 92.28/36.59 & \textbf{93.50}/23.58 & 92.28/34.96\\
    100 & 92.68/30.49 & 93.09/27.24 & \textbf{93.50}/15.85 & 91.46/51.22\\
    200 & 92.68/39.43 & 92.68/32.11 & 93.09/41.46 & \textbf{93.50}/27.24\\
    \bottomrule
  \end{tabular}
\end{table}

\subsection{RQ5: Computer Science QA}

Computer Science is a harder routing environment. All-Mini answers 80.49\% of
the 205 test questions correctly at a realized cost of \$0.23048; all-Nano
answers 71.22\% at \$0.05960. The outcome oracle reaches 83.90\% at \$0.08243,
so complementary Nano-only successes again create theoretical headroom.

The strongest accuracy-first learned result is a five-fold out-of-fold soft
prompt ensemble. Its global threshold exactly matches all-Mini accuracy, but
routes 200 of 205 questions to Mini and saves only 1.29\% of answer-model cost.
A local Qwen3.5-2B router obtains 80.00\% accuracy and 4.26\% saving, missing one
of 26 critical cases. Seeding GEPA with Mini-biased demonstrations creates a
more aggressive policy: it sends 140 questions to Mini, reaches 78.54\%, and
saves 17.65\% of cost. The non-Mini-biased GEPA run collapses to all-Nano.
These results mirror Movie: prompt optimization can move the operating point,
but scarce critical outcomes make preserving the all-Mini target expensive.

\begin{table}[t]
  \centering
  \small
  \caption{MMLU-Pro Computer Science fixed test (205 questions). Cost saving is
  relative to the \$0.23048 all-Mini answer-model cost.}
  \label{tab:qa}
  \begin{tabular}{lrrr}
    \toprule
    Method & Accuracy & Mini rate & Cost saving\\
    \midrule
    All Nano & 71.22\% & 0.00\% & 74.14\%\\
    All Mini & 80.49\% & 100.00\% & 0.00\%\\
    Outcome oracle & 83.90\% & 12.68\% & 64.24\%\\
    OOF soft-prompt ensemble & 80.49\% & 97.56\% & 1.29\%\\
    Local Qwen hard router & 80.00\% & 93.66\% & 4.26\%\\
    GEPA, Mini-biased seed & 78.54\% & 68.29\% & 17.65\%\\
    GEPA, default seed & 71.22\% & 0.00\% & 74.14\%\\
    \bottomrule
  \end{tabular}
\end{table}

\subsection{Interim Answers and Remaining Experiments}

The present evidence answers RQ1 cautiously: historical savings are real on the
recorded split, but no direct router both matches all-Mini and produces a useful,
stable fresh saving. For RQ2, soft prompts dominate the current discrete GEPA
budget on historical accuracy, while both fail under rare-event shift. RQ3 shows that, after the same full paired screen, disagreement-guided
selection uses a fixed human-label budget to recover substantially more
one-model-only cases than random annotation. The current downstream router does
not yet exploit all of that coverage; a matched
inverse-expected-cost WCE rerun is still required. RQ4 rejects a
monotonic ``more tokens or labels is better'' hypothesis. RQ5 shows that the
fresh and second-task results are materially less optimistic than historical
Movie selection.

\todo{Before submission, execute the pre-registered final matrix: (1) train the lightweight soft-prompt with inverse-expected-cost WCE and rerun
the discrete method on identical splits under the shared cost score; (2)
draw a new 300--500 review manifest from the remaining unpaired SemBench IDs;
(3) run GEPA hard and scored routing for three seeds and at least three
metric-call budgets; (4) add matched-Mini-budget random, FrugalGPT-style
confidence, and RouteLLM baselines; (5) repeat on at least two additional
semantic operators, including one join or map; (6) measure router latency and
dollar overhead; and (7) report paired bootstrap confidence intervals, McNemar
tests against all-Mini, and sensitivity to price snapshots and utility
normalization. The already-consumed 200-row set must not be relabeled as fresh.}
